% ===== File: proposal_rq2.tex =====
\documentclass[sigplan,nonacm,10pt]{acmart}
\settopmatter{printfolios=true}

\usepackage{booktabs}
\usepackage{enumitem}
\usepackage{amsmath}
\setlist{nosep,leftmargin=*}

\begin{document}

\title{Predicting \texttt{cp.async} Pipeline Effectiveness\\
from Memory-Subsystem Parameters: A Cross-GPU Study}

\author{Xiangchen Song}
\affiliation{\institution{University of Michigan -- Ann Arbor}}
\email{xxxchen@umich.edu}
\author{Chenxin Geng}
\affiliation{\institution{University of Michigan -- Ann Arbor}}
\email{gchenxin@umich.edu}
\author{Yuyang Feng}
\affiliation{\institution{University of Michigan -- Ann Arbor}}
\email{yuyangfe@umich.edu}
\author{Tian Xia}
\affiliation{\institution{University of Michigan -- Ann Arbor}}
\email{crazyjoe@umich.edu}

% ────────────────────────────────────────────────────────────
\begin{abstract}
% ────────────────────────────────────────────────────────────
NVIDIA's \texttt{cp.async} instruction enables asynchronous
global-to-shared memory copies that can be software-pipelined with
computation, and is now standard in high-performance GPU kernels.
Practitioners typically tune pipeline depth per kernel and assume the
result transfers across GPUs sharing the same instruction.
We test this assumption by running \emph{identical} \texttt{cp.async}
pipelined kernels---FP32 GEMM and a 1D stencil---on three GPUs that
support \texttt{cp.async} but have radically different memory subsystems:
A40 (GDDR6, 6\,MB L2, 696\,GB/s, ridge\,=\,53.7),
A100 MIG~3g.40gb (HBM2e, 20\,MiB L2, 968\,GB/s, ridge\,=\,7.9), and
H100 SXM5 (HBM3, 50\,MB L2, 3350\,GB/s, ridge\,=\,20.0).
We build a parameter-free predictive model from architecture constants
and microbenchmarked latencies, validate it across all three GPUs and
both kernels, and package it as a lightweight Triton autotune pre-filter
that prunes the pipeline-depth search space.
A pointer-chasing microbenchmark characterizes each GPU's effective L2
behavior---including an empirical characterization of MIG-instance L2
capacity and latency, which is not specified in vendor
documentation---and feeds directly into the predictor.
\end{abstract}

\maketitle

% ============================================================
\section{Problem Definition and Research Question}
\label{sec:problem}
% ============================================================

Hiding global memory latency is central to GPU kernel optimization.
NVIDIA's \texttt{cp.async} (Ampere and later) enables asynchronous
global-to-shared copies that bypass registers and can be
software-pipelined with computation~\cite{cuda_guide}.
Libraries such as CUTLASS~3.x~\cite{cutlass3} and
FlashAttention-3~\cite{flashattn3} rely heavily on multi-stage
\texttt{cp.async} pipelines, and autotuners (including
Triton's~\cite{triton}) sweep pipeline depth $S$ as a key tuning knob.

However, pipeline effectiveness is not purely an ISA property.
It depends on the interaction between L2 cache capacity (determining
whether tile data is served from cache or DRAM), memory bandwidth and
latency (determining how long a copy stalls), and per-SM shared memory
(determining how many pipeline stages fit before occupancy degrades).
These parameters vary dramatically even among GPUs that share
\texttt{cp.async} support---yet no study has quantified how this
variation affects optimal pipeline configuration, and no existing
autotuner uses architecture parameters to prune the search space.

\smallskip
\noindent\textbf{Research Question:}
\emph{For an identical \texttt{cp.async} pipeline kernel, does the optimal
pipeline depth $S^*$ and its marginal speedup diverge across GPUs with
different memory subsystems?  Can a lightweight, parameter-free model
predict this divergence from architecture constants alone---and can such
a model reduce autotuning cost in practice?}

% ============================================================
\section{Hypotheses}
\label{sec:hypotheses}
% ============================================================

We test three competing explanations for how memory-subsystem parameters
govern \texttt{cp.async} pipeline effectiveness.  Each hypothesis makes
distinct, falsifiable predictions about the observables defined in
Section~\ref{sec:plots}.

\begin{itemize}
\item \textbf{H1 (L2-capacity-driven):}
  Pipeline speedup exhibits a sharp inflection when the concurrent
  working set $W_{\text{conc}}$ exceeds L2 capacity $C_{\text{L2}}$.
  Plotting speedup against $W_{\text{conc}}/C_{\text{L2}}$ collapses
  all three GPU curves onto a single master curve.
  %
  \emph{Falsifiable prediction:} after L2 normalization, A100-MIG
  (20\,MiB) and H100 (50\,MB) should \emph{not} overlap---if they do
  despite a 2.4$\times$ L2 gap, L2 capacity alone is insufficient.

\item \textbf{H2 (Overlap-saturation / ridge-point-driven):}
  Pipeline speedup is governed by whether per-tile compute time can mask
  copy latency, i.e.\ by the ridge point
  $R = F_{\text{peak}}/\text{BW}_{\text{peak}}$.
  GPUs with low $R$ (more memory-bound) benefit most from pipelining.
  %
  \emph{Falsifiable prediction:} A100-MIG ($R = 7.9$, lowest of the
  three) should show the largest pipeline speedup; A40 ($R = 53.7$)
  the smallest---possibly even a slowdown if occupancy loss outweighs
  negligible latency-hiding benefit.

\item \textbf{H3 (Latency--occupancy tension):}
  Optimal depth $S^*$ results from the tension between memory latency
  (favoring deeper pipelines) and shared-memory capacity (limiting depth
  through occupancy loss).
  A40 has the highest expected DRAM latency (GDDR6; latency to be
  confirmed by microbenchmark, Section~\ref{sec:mig_microbench}) yet
  the smallest per-SM shared memory (100\,KB).
  %
  \emph{Falsifiable prediction:} $S^*(N)$ is non-monotonic on A40
  (rises then falls with $N$), while remaining monotone or flat on
  A100-MIG and H100 (more SMEM headroom).
\end{itemize}

\noindent\textbf{Discrimination logic.}
H1: curve-collapse under L2 normalization.
H2: speedup rank follows ridge-point rank, not L2 rank.
H3: $S^*(N)$ exhibits GPU-dependent non-monotonicity.
In practice a mixture is likely; the predictive model
(Section~\ref{sec:model}) incorporates all three mechanisms and
the \emph{pipeline benefit map} (Section~\ref{sec:heatmap})
visualizes their interplay across the full (AI, $W/C_{\text{L2}}$)
space.

% ============================================================
\section{Related Work}
\label{sec:related}
% ============================================================

\textbf{Roofline model.}
Williams et al.~\cite{roofline} introduced the roofline model relating
arithmetic intensity to attainable performance.  We use it as the
primary cross-GPU normalization framework, computing per-GPU roofline
ceilings from runtime-queried hardware parameters (including
MIG-instance peaks).

\textbf{\texttt{cp.async} pipeline studies.}
Li et al.~\cite{dtcspmm} demonstrate $>$2$\times$ speedup from
double-buffered \texttt{cp.async} pipelines for SpMM on Ampere (ASPLOS
2024); our V3 adopts a similar multi-stage pattern but sweeps stage
count across three GPUs and two kernels rather than one.
Chen et al.~\cite{tawa} (Tawa, 2025) sweep pipeline depth on H100,
revealing occupancy--depth tradeoffs that motivate H3; we test whether
these tradeoffs shift on GPUs with smaller SMEM.
Wu et al.~\cite{twill} (Twill, 2025) note that optimal pipelining
relies on ``brittle heuristics''; our Triton pre-filter
(Section~\ref{sec:triton}) targets a complementary, spec-based
alternative.

\textbf{GPU L2 cache and tile scheduling.}
Wei et al.~\cite{hytis} (HyTiS, SC 2025) model L2-aware tile
scheduling for GEMM, showing up to 36\% throughput variation from tile
placement on A100/H100, but do not study async pipeline interaction
at L2 boundaries.
Jin et al.~\cite{gpu_noc} (MICRO 2024) reverse-engineer GPU NoC
topology and find up to 70\% L2 latency variation from SM-to-slice
distance.
Ernst et al.~\cite{layer_conditions} identify GPU cache ``layer
conditions'' where performance shifts qualitatively at capacity
thresholds.
No prior work characterizes L2 cache behavior under MIG partitioning.

\textbf{GPU microbenchmarking.}
Mei and Chu~\cite{mei_dissecting} establish pointer-chasing methodology
for GPU cache hierarchy characterization.
Wang et al.~\cite{hopper_microbench} microbenchmark Hopper's TMA and
L2 on H100 only.
We extend pointer-chasing to MIG instances, measuring effective
(not nominal) L2 capacity and latency under partitioned resources.

\textbf{Cross-GPU benchmarking.}
Svedin et al.~\cite{svedin_lineage} benchmark K80 through A100 on
Rodinia (2021) but predate H100, lack roofline framing, and omit A40.
No prior work compares the \emph{same} \texttt{cp.async} pipeline
kernel across GPUs with different memory subsystems.

\textbf{Analytical models and autotuning.}
Guo et al.~\cite{tritonblas} (tritonBLAS, 2025) propose an analytical
kernel-parameter selector for Triton GEMM but target a single GPU.
ACTA~\cite{acta} models TMA configuration on Hopper.
Triton's \texttt{@triton.autotune}~\cite{triton} performs exhaustive
grid search over all configurations including \texttt{num\_stages}.
Our model is simpler (overlap + occupancy penalty) but designed to
\emph{transfer across GPUs} and to \emph{prune the autotune search
space} before exhaustive search begins.

% ============================================================
\section{Experimental Setup}
\label{sec:setup}
% ============================================================

\subsection{Hardware and Experimental Roles}

\begin{table}[t]
\centering
\small
\caption{Target GPUs.  A100 values are runtime
\texttt{cudaDeviceProp} measurements from the MIG~3g.40gb instance.
All other GPUs are full devices.}
\label{tab:gpus}
\begin{tabular}{@{}lcccc@{}}
\toprule
 & \textbf{V100} & \textbf{A40} & \textbf{A100 MIG 3g} & \textbf{H100} \\
\midrule
Cluster        & Great Lakes & Great Lakes & Great Lakes & Lighthouse \\
Arch / SM ver. & Volta / 70  & Ampere / 86 & Ampere / 80 & Hopper / 90 \\
SMs            & 80          & 84          & 42          & 132 \\
\texttt{cp.async} & ---      & \checkmark  & \checkmark  & \checkmark \\
\midrule
Memory tech    & HBM2        & GDDR6       & HBM2e (MIG) & HBM3${}^\dagger$ \\
GPU memory     & 32\,GB      & 48\,GB      & 40\,GB      & 80\,GB \\
Bus width      & 4096\,bit   & 384\,bit    & 2560\,bit   & 5120\,bit \\
Peak BW        & 900\,GB/s   & 696\,GB/s   & 968\,GB/s   & 3350\,GB/s \\
\midrule
L2 cache       & 6\,MB       & 6\,MB       & 20\,MiB${}^\ddagger$  & 50\,MB \\
SMEM/SM        & 96\,KB      & 100\,KB     & 164\,KB     & 228\,KB \\
\midrule
FP32 peak      & 15.7\,TF    & 37.4\,TF    & 7.6\,TF     & 67\,TF \\
Ridge (FLOP/B) & 17.4        & 53.7        & 7.9         & 20.0 \\
\midrule
\textbf{Role}  & Boundary    & Core anchor & Intermediate & Core anchor \\
\bottomrule
\multicolumn{5}{@{}l@{}}{\footnotesize ${}^\dagger$H100 on Lighthouse is the SXM5 variant (HBM3, 3.35\,TB/s).} \\
\multicolumn{5}{@{}l@{}}{\footnotesize ${}^\ddagger$Reported by \texttt{cudaDeviceProp}; effective capacity $C_{\text{L2}}^{\text{eff}}$ measured via microbenchmark.}
\end{tabular}
\end{table}

The four GPUs (Table~\ref{tab:gpus}) play distinct roles:

\begin{itemize}
\item \textbf{A40 and H100 (core contrast pair, both full GPUs):}
  Maximum memory-subsystem divergence among \texttt{cp.async}-capable
  GPUs: 8.3$\times$ L2, 4.8$\times$ BW, 2.7$\times$ ridge point---with
  no MIG artifacts.

\item \textbf{A100 MIG 3g.40gb (intermediate + ridge-point extreme):}
  MIG partitions compute more aggressively than bandwidth
  (42/108\,=\,39\% SMs, 50\% memory bus), producing the \emph{lowest}
  ridge point (7.9) of any GPU in the study.
  H2 predicts it should benefit \emph{most} from pipelining.
  Also serves as a real-world test of predictor robustness under
  MIG-partitioned resources~\cite{mig_guide}.

\item \textbf{V100 (boundary condition):}
  Lacks \texttt{cp.async}; runs V0--V1 only.  The predictor must output
  speedup\,=\,1.0$\times$.
\end{itemize}

\smallskip
\noindent\textbf{GPU pair separation logic.}
The three \texttt{cp.async}-capable GPUs span a wide, non-aligned
region of (L2, BW, ridge) space:

{\small
\begin{center}
\begin{tabular}{@{}lcccl@{}}
\toprule
\textbf{Pair} & $\Delta$\textbf{L2} & $\Delta$\textbf{BW}
  & $\Delta$\textbf{Ridge} & \textbf{Primary contrast} \\
\midrule
A40 vs.\ H100       & 8.3$\times$ & 4.8$\times$ & 2.7$\times$
  & Max.\ memory-subsystem gap \\
A40 vs.\ A100-MIG   & 3.3$\times$ & 1.4$\times$ & 6.8$\times$
  & Same ISA; ridge-point extreme \\
A100-MIG vs.\ H100  & 2.4$\times$ & 3.5$\times$ & 2.5$\times$
  & Intermediate $\to$ high \\
\bottomrule
\end{tabular}
\end{center}
}

\noindent No two parameters are monotonically ordered across all three
GPUs (A100-MIG has the \emph{lowest} ridge despite mid-range L2), so
the predictor cannot succeed by fitting a single monotonic trend.

\subsection{Kernel Variants}

\begin{table}[t]
\centering
\small
\caption{Kernel variants.  V3-GEMM and V3-Stencil are the
\texttt{cp.async} pipeline kernels central to the study.}
\label{tab:variants}
\begin{tabular}{@{}llp{5.2cm}@{}}
\toprule
\textbf{ID} & \textbf{GPUs} & \textbf{Definition} \\
\midrule
V0-GEMM  & All   & Naive global load, no tiling. \\
V1-GEMM  & All   & Shared-memory tiling, synchronous load +
                    \texttt{\_\_syncthreads}. \\
V3-GEMM  & A40+  & Multi-stage \texttt{cp.async} pipeline,
                    $S \in \{2, 3, 4\}$.  Tile geometry matches V1. \\
\midrule
V1-Stencil & All  & 1D stencil, synchronous shared-memory tiling
                     with halo. \\
V3-Stencil & A40+ & Same stencil with \texttt{cp.async} pipelined
                     loads, $S \in \{2, 3, 4\}$. \\
\midrule
Ref & All   & cuBLAS SGEMM (ceiling / sanity). \\
\bottomrule
\end{tabular}
\end{table}

We implement two kernel families (Table~\ref{tab:variants}):

\textbf{FP32 GEMM} ($C = A \times B$, row-major) with matched tile
structures across V0, V1, and V3.
GEMM has high arithmetic intensity ($\text{AI} \gg R$ for all GPUs at
$N \ge 512$), making it compute-bound in the roofline sense but still
subject to per-tile memory stalls that pipelining can hide.

\textbf{1D Stencil} with a 5-point stencil on a 1D array of length $L$.
Stencil has much lower AI ($\sim$5 FLOP/B after tiling), placing it
near or below the ridge point of most GPUs.  This makes it genuinely
memory-bound and provides a complementary test:
if the predictor works for both a compute-bound kernel (GEMM) and a
memory-bound kernel (stencil), the model captures general
memory-subsystem properties rather than GEMM-specific behavior.
To ensure the pipeline has sufficient compute to overlap with async
copies (rather than being dominated by barrier overhead), the stencil
kernel uses loop unrolling and multiple elements per thread to increase
per-stage compute density.

The \textbf{same source code} for V3-GEMM (and V3-Stencil) compiles and
runs on A40, A100-MIG, and H100 with no GPU-specific branches.

\subsection{Workload and Measurement}

\textbf{GEMM:} $N \in \{256, 512, 1024, 2048, 3072, 4096, 6144, 8192\}$.
For the \emph{pipeline benefit map} (Section~\ref{sec:heatmap}), we
additionally sweep tile sizes
$T \in \{16, 32, 64, 128\}$ to vary per-tile AI.

\textbf{Stencil:}
$L \in \{2^{16}, 2^{18}, 2^{20}, 2^{22}, 2^{24}, 2^{26}\}$.

Timing: \texttt{cudaEvent\_t}, 3 warmup + 10 timed iterations, median.
Correctness: verified against reference (cuBLAS for GEMM; CPU for stencil).

% ============================================================
\section{MIG L2 Cache Microbenchmark}
\label{sec:mig_microbench}
% ============================================================

NVIDIA documents that MIG partitions memory by slice, but does not
publish effective L2 capacity, associativity, or replacement behavior
under MIG~\cite{mig_guide}.
The \texttt{cudaDeviceProp}-reported 20\,MiB may not reflect the
\emph{usable} capacity if MIG boundaries cause additional conflict
misses or if L2 partition granularity differs from the nominal slice
assignment.

We run a \textbf{pointer-chasing microbenchmark}
(following Mei and Chu~\cite{mei_dissecting}) on every GPU to measure:

\begin{enumerate}
\item \textbf{L2 effective capacity ($C_{\text{L2}}^{\text{eff}}$):}
  Sweep array size from 1\,MB to 40\,MB; identify the inflection where
  average load latency jumps from L2-resident to DRAM-resident.
  On A100-MIG, if $C_{\text{L2}}^{\text{eff}} < 20$\,MiB, MIG imposes
  hidden overhead.

\item \textbf{L2 hit latency ($L_{\text{L2}}$) and DRAM latency
  ($L_{\text{DRAM}}$):}
  Measured in cycles at plateau regions below and above
  $C_{\text{L2}}^{\text{eff}}$.
  These are direct inputs to the predictive model
  (Section~\ref{sec:model}).

\item \textbf{L2 bandwidth curve:}
  Streaming bandwidth vs.\ working-set size, per GPU.
  Reveals whether MIG's L2 partition sustains the same bandwidth per
  byte as the full die.
\end{enumerate}

These measurements are treated as \emph{per-GPU constants} (not
fitted from GEMM data) and are performed once during Week~1--2.
If $C_{\text{L2}}^{\text{eff}}$ differs from the reported 20\,MiB on
the MIG instance, this provides an empirical characterization of
MIG-instance effective L2 behavior---information not specified in
vendor documentation~\cite{mig_guide}.

% ============================================================
\section{Working-Set Definition and L2 Analysis}
\label{sec:ws}
% ============================================================

For a tiled GEMM with tile dimensions $T_M \times T_K$ and
$T_K \times T_N$, the per-block working set is:
\[
  W_{\text{block}} = (T_M \cdot T_K + T_K \cdot T_N) \times 4
  ~\text{bytes (FP32)}.
\]
The concurrent working set across simultaneously active blocks sharing
the L2:
\[
  W_{\text{conc}} = W_{\text{block}} \times S \times
  \min\!\bigl(B_{\text{active}},\;
  \lceil N/T_M \rceil \cdot \lceil N/T_N \rceil\bigr),
\]
where $S$ is the pipeline stage count and
$B_{\text{active}} = N_{\text{SM}} \times
\lfloor S_{\text{SM}} / (W_{\text{block}} \cdot S) \rfloor$
is the occupancy-limited concurrency.

For the stencil, $W_{\text{block}}$ includes the tile plus halo
elements; the same $W_{\text{conc}}$ framework applies with adjusted
tile geometry.

In both cases, we use the \emph{microbenchmarked}
$C_{\text{L2}}^{\text{eff}}$ (Section~\ref{sec:mig_microbench}) rather
than the nominal \texttt{cudaDeviceProp} value when computing
$W_{\text{conc}} / C_{\text{L2}}$.

% ============================================================
\section{Key Plots and Hypothesis Discrimination}
\label{sec:plots}
% ============================================================

\begin{enumerate}
\item \textbf{Speedup vs.\ $N$ (raw), per kernel:}
  Pipeline speedup for A40, A100-MIG, H100 on the same plot.
  Separate plots for GEMM and stencil.
  If all three curves separate, ``pipeline tuning transfers'' is
  falsified.

\item \textbf{Speedup vs.\ $W_{\text{conc}}/C_{\text{L2}}^{\text{eff}}$
  (L2-normalized):}
  Tests H1 (collapse $\Rightarrow$ L2 explains all).
  Uses microbenchmarked effective L2, not nominal.

\item \textbf{$S^*(N)$ per GPU:}
  Optimal pipeline depth as a function of problem size.
  Tests H3 (non-monotonicity on A40).

\item \textbf{Speedup rank vs.\ ridge-point rank:}
  Tests H2 (A100-MIG $>$ H100 $>$ A40 ordering).

\item \textbf{Model prediction vs.\ measurement (scatter):}
  $x$ = predicted, $y$ = measured, one point per $(N, g, S)$ triple.
  MAPE per GPU and per kernel.
\end{enumerate}

\subsection{Pipeline Benefit Map}
\label{sec:heatmap}

Beyond the above plots, we produce a \textbf{two-dimensional heatmap}
for each GPU:
\begin{itemize}
\item $x$-axis: per-tile arithmetic intensity (varied by sweeping tile
  size $T \in \{16, 32, 64, 128\}$ in GEMM);
\item $y$-axis: $W_{\text{conc}} / C_{\text{L2}}^{\text{eff}}$
  (varied by sweeping $N$);
\item color: pipeline speedup $T_{V1}/T_{V3_{\text{best}}}$.
\end{itemize}

This map answers the practical question: \emph{for a kernel with a
given AI and data footprint, on this GPU, is pipeline worth enabling?}
The shape of the ``pipeline-beneficial'' region (green) differs across
GPUs:
on A40 (ridge\,=\,53.7) only the low-AI, high-$W/C$ corner should be
green;
on A100-MIG (ridge\,=\,7.9) most of the map should be green.
Comparing the three maps provides a visual summary of all three
hypotheses simultaneously.

\subsection{Profiling Signals}

Nsight Compute metrics collected at each $(N, g, S)$:
L2 hit rate, DRAM bytes read, achieved occupancy, barrier stall cycles,
and L2 throughput utilization.
These provide mechanistic support: a speedup inflection aligned with
L2-hit-rate drop supports H1; speedup tracking compute-to-copy ratio
(barrier stalls) supports H2; non-monotonic $S^*$ reversing when
occupancy drops supports H3.

\subsection{Anomaly Analysis Protocol}
\label{sec:anomaly}

After collecting all data, we compute
$\Delta(N,g,S) = \widehat{\text{Speedup}} - \text{Speedup}_{\text{measured}}$
for every configuration and rank by $|\Delta|$.
The top-5 largest deviations receive a targeted Nsight Compute
deep-dive: full metric collection at that specific $(N,g,S)$ to
identify the root cause of model error.

Possible anomaly sources include:
MIG L2 partitioning artifacts (addressed by
Section~\ref{sec:mig_microbench});
L2 bank conflicts at specific tile geometries (cf.~\cite{gpu_noc});
occupancy cliff effects where one additional SMEM buffer crosses a
block-per-SM boundary.
\textbf{Documenting model failure modes is itself a contribution}---it
identifies which architectural effects a simple overlap model misses
and where more detailed modeling is needed.

% ============================================================
\section{Predictive Model (Parameter-Free)}
\label{sec:model}
% ============================================================

We construct a semi-analytic predictor of \texttt{cp.async} pipeline
speedup combining an L2-aware latency estimate with an occupancy
penalty, following the overlap reasoning of Little's Law
(cf.~\cite{acta}).

\subsection{Effective Memory Latency}

Let $h$ be the estimated L2 hit fraction:
\[
  h(N,g) = \min\!\bigl(1,\;
  C_{\text{L2}}^{\text{eff}}(g) \;/\; W_{\text{conc}}(N,g)\bigr).
\]
Note: $C_{\text{L2}}^{\text{eff}}$ is the \emph{microbenchmarked}
effective capacity (Section~\ref{sec:mig_microbench}), not the nominal
device-property value.
Our $h$ is a capacity-based upper bound on the true hit rate; real L2
behavior is further affected by set-associativity, slice locality, and
replacement policy~\cite{gpu_noc}.
Deviations between this approximation and reality are explicitly
analyzed via the anomaly protocol (Section~\ref{sec:anomaly}) and
Nsight Compute L2 hit-rate metrics.
The effective latency seen by a tile load is:
\[
  L_{\text{eff}}(N,g) = h \cdot L_{\text{L2}}(g)
  + (1 - h) \cdot L_{\text{DRAM}}(g),
\]
where $L_{\text{L2}}$ and $L_{\text{DRAM}}$ are microbenchmarked
constants.

\subsection{Predicted Speedup}

Minimum depth to fully hide latency:
\[
  S_{\min}(N,g) = \left\lceil
  \frac{L_{\text{eff}}(N,g)}{T_{\text{compute}}(g)}
  \right\rceil,
\]
where $T_{\text{compute}}$ is per-tile pure compute time, isolated from
memory stalls via one of two methods:
(i)~Nsight Compute's \texttt{smsp\_\_cycles\_active} metric on V1
minus memory-stall cycles, or
(ii)~a compute-only microkernel that operates on data already resident
in shared memory / registers.
Using raw V1 wall-clock time would conflate memory stalls with compute,
biasing $S_{\min}$; isolating pure compute avoids this circularity.
Predicted speedup:
\[
  \widehat{\text{Speedup}}(N, g, S) =
  \frac{\max\bigl(T_{\text{compute}},\; L_{\text{eff}}\bigr)}
  {\max\bigl(T_{\text{compute}},\;
   L_{\text{eff}} / \min(S,\, S_{\min})\bigr)}
  \;\times\;
  \frac{\text{Occ}(S,g)}{\text{Occ}(1,g)}.
\]
The occupancy ratio $\text{Occ}(S)/\text{Occ}(1)$ is computed from the
CUDA occupancy API given measured SMEM and register usage---not fitted.

\subsection{Why ``Parameter-Free'' Despite Few GPUs}

Every model input is either an architecture constant
($C_{\text{L2}}^{\text{eff}}$, BW, $F_{\text{peak}}$, $S_{\text{SM}}$,
$N_{\text{SM}}$) or a per-GPU microbenchmark constant ($L_{\text{L2}}$,
$L_{\text{DRAM}}$).
No free parameters are fitted from GEMM or stencil results.
The kernel experiments serve to \emph{validate}, not fit.
Demonstrating accuracy on two distinct kernels (GEMM and stencil)
across three GPUs further guards against overfitting to a single
workload.

\subsection{Validation Protocol}

\begin{enumerate}
\item \textbf{Per-GPU accuracy:} predicted vs.\ measured speedup for
  all $(N, S)$ on each GPU.  Report MAPE per GPU, per kernel.
\item \textbf{Boundary condition:} V100 $\to$ 1.0$\times$ by
  construction.
\item \textbf{Leave-one-out:} hide one GPU's \emph{measurements},
  predict from its architectural parameters alone, report held-out MAPE.
\item \textbf{Nominal vs.\ effective L2:} compare MAPE when using
  $C_{\text{L2}}^{\text{eff}}$ (microbenchmarked) vs.\
  $C_{\text{L2}}^{\text{nom}}$ (\texttt{cudaDeviceProp}).
  If microbenchmarked values improve accuracy, the L2 characterization
  has direct predictive value.
\end{enumerate}

% ============================================================
\section{Triton Autotune Pre-Filter}
\label{sec:triton}
% ============================================================

We package the predictive model as a \textbf{lightweight pre-filter for
Triton's \texttt{@triton.autotune} decorator}~\cite{triton}.

\textbf{Current Triton behavior.}
Triton's autotune exhaustively benchmarks every configuration in the
user-specified grid, including all \texttt{num\_stages} values.
For a typical GEMM tuning grid with 4 tile sizes $\times$ 3 stage
counts $\times$ 2 \texttt{num\_warps} settings = 24 configurations,
each evaluated with warmup + timed iterations, autotuning can take
30--60 seconds per kernel launch shape.

\textbf{Our pre-filter.}
Before autotune begins, query the GPU's architectural parameters
(via \texttt{cudaDeviceProp} and our microbenchmarked latencies),
compute the predicted speedup for each \texttt{num\_stages} candidate,
and \emph{prune} candidates whose predicted speedup is within a
threshold $\epsilon$ of 1.0$\times$ (i.e., pipeline unlikely to help).

\textbf{Expected outcome.}
On A40 (ridge\,=\,53.7), the pre-filter should prune $S \in \{3, 4\}$
for large $N$ (occupancy cost exceeds benefit), reducing the search
space by $\sim$50\%.
On A100-MIG (ridge\,=\,7.9), few or no candidates are pruned (pipeline
is beneficial).
We measure:
(i)~wall-clock autotune time with and without pre-filter,
(ii)~whether pruning ever removes the true optimal configuration
(false-negative rate), and
(iii)~end-to-end kernel performance (must not degrade).

\textbf{Implementation.}
$\sim$150 lines of Python wrapping the model equations from
Section~\ref{sec:model}, invoked as a decorator that wraps
\texttt{@triton.autotune}.
Evaluated on Triton's tutorial GEMM kernel and (if time permits) on
a Triton stencil kernel.

% ============================================================
\section{Milestones, Schedule, and Division of Labor}
\label{sec:plan}
% ============================================================

\begin{table}[t]
\centering
\small
\caption{Timeline and responsibilities.  Items marked ${}^\star$ are
stretch goals that enhance novelty but are not required for the core
contribution.}
\label{tab:plan}
\begin{tabular}{@{}cp{5.2cm}l@{}}
\toprule
\textbf{Week} & \textbf{Milestone} & \textbf{Lead} \\
\midrule
1--2 & Build system (CMake, SLURM).
       V0/V1 baselines on all 4 GPUs.
       MIG device-query validation.
       \textbf{Pointer-chasing microbench}
       ($L_{\text{L2}}$, $L_{\text{DRAM}}$,
       $C_{\text{L2}}^{\text{eff}}$) on all GPUs.
     & A + C \\
3--4 & V3-GEMM ($S = 2,3,4$) + V1/V3-Stencil
       implementation.
       Initial speedup curves on A40, A100-MIG.
     & B \\
5--6 & Full A40 / A100-MIG / H100 GEMM comparison.
       Nsight profiling.
       Tile-size sweep for heatmap data.
       \textit{Milestone~1 slides.}
     & B + D \\
7--8 & Stencil experiments on all GPUs.
       Predictor implementation + validation.
       Anomaly deep-dives.
       ${}^\star$Triton pre-filter prototype.
     & C + D \\
9--10 & Pipeline benefit maps.
       Final figures, hypothesis conclusions.
       Report + poster.
     & All \\
\bottomrule
\end{tabular}
\end{table}

\noindent\textbf{Member roles:}
\textbf{A}~build system, baselines, MIG validation, pointer-chasing
microbench;
\textbf{B}~V3-GEMM and V3-Stencil implementation, A40/A100-MIG runs,
tile-size sweep;
\textbf{C}~H100 coordination (Lighthouse), stencil runs,
Triton pre-filter;
\textbf{D}~Nsight profiling, L2/working-set analysis, model
implementation, heatmap generation, anomaly analysis.
All members contribute to writing and poster.

\smallskip
\noindent\textbf{Decision point (Week~6).}
If GEMM speedups converge across GPUs (no divergence), the
``pipeline tuning is transferable'' conclusion is itself useful; we
then invest more effort in the stencil kernel (lower AI, higher chance
of divergence) and the MIG L2 characterization.
If divergence is strong, we prioritize the predictor, heatmap, and
Triton integration.

\smallskip
\noindent\textbf{Contributions at a glance.}
\begin{enumerate}
\item A controlled cross-GPU \texttt{cp.async} pipeline comparison
      (same code, three memory subsystems, two kernels).
\item An empirical characterization of MIG-instance effective L2
      behavior (capacity and latency under pointer-chasing), not
      specified in vendor documentation.
\item A parameter-free predictor validated on two kernels and three GPUs
      (including MIG).
\item A pipeline benefit map visualizing the (AI, $W/C_{\text{L2}}$)
      region where pipelining helps, per GPU.
\item A Triton autotune pre-filter prototype demonstrating practical
      search-space reduction.
\end{enumerate}

% ============================================================
\bibliographystyle{ACM-Reference-Format}
\bibliography{refs}

\end{document}
